\chapter{Examples}
\emph{In this chapter we demonstrate a few simple examples of how our embedding of Lilac in Iris can be used for interactively proving specifications of APPL programs.
We develop some proof automation, and reflect on more divergences from the Lilac paper \cite{lilac}}


First in sec \ref{sec:uniform} we verify the most basic programs which we had struggled to verify under the original definitions in Chapter \ref{sec:definitions}. Next
in sec \ref{sec:expectation} we outline difficulties faced with reasoning about \texttt{LProp}s which used the \texttt{expectation} ($\mathbb{E}[X]=e$) connective in the logic.
We develop some proof rules regarding expectation and some custom tactics for ergonomic reasoning.

Most notably we use this example to remark on how in our mechanised proofs we completely sidestep the usage of almost surely equals connective ($\circeq$) and substitutions using the congruence rule.
This is not specific to this particular example and is a general phenomenon. Indeed, the Lilac paper \cite{lilac} has not provided proof rules for introduction of
$\circeq$\footnote{the original notation from the paper for the almost surely equals connective is $\overset{as}{=}$. We use $\circeq$ in the lean formalisation}, and uses it rather intuitively
in the proof sketches. In our formalisation, we have required it at all.

\section{Sampling Uniform Distributions}
\label{sec:uniform}

We recall the proof from Sec \ref{sec:parametrised}, which we had failed to complete under the old definitions. 
The informal and lean version, taken from \repo{Examples}, are provided side-by-side.

\begin{minipage}[t]{0.48\textwidth}
\begin{leancode}
abbrev unif1 : Term [] Ty.real.G :=
  unif01.bind (
  ret (#0)
)
\end{leancode}
\end{minipage}
\hfill
\begin{minipage}[t]{0.48\textwidth}
\[
x \leftarrow \mathrm{unif}\ [0,1];\\
 \mathrm{ret}\ x
\]
\end{minipage}

Since we are using de Brujin indices we cannot name variables bound by the \texttt{bind} construct, but instead need to access this variables
from an implicit context. \texttt{\#0} is shorthand for \texttt{Term.var Member.head}, which refers to the syntactically closest (right most)
variable bound. In this case it is obvious which variable we refer to since there is only a single variable. Below the informal
specification is followed by the formalised Lean specification:

\[
\{ \top \}\ \textsc{Unif1}\ \{ X.\ X \sim \mathrm{Unif}[0,1] \}
\]

\begin{leancode}
nil : RV (List.TProd (⟪·⟫) []) := ...
abbrev post1 : LProp :=
  wp (unif1.den ∘ᵣ nil) (λ X : RV ⟪Ty.real⟫ ↦ (X ∼ unif01_sem))
\end{leancode}

We see that the lean code of the specification has some additional components implicit in the informal version.  Recall that \texttt{wp} has type, \lean{def wp (M : RV (⟪Ty.G A⟫)) (Q : RV ⟪A⟫ → LProp) : LProp}.
We have assigned \lean{(M := unif1.den ∘ᵣ nil)}. Why is it not just \texttt{unif1.den}, the denotation of the unif1 program we had just defined above?
Inspecting the types, $\text{unif1.den} : \sem{\Delta} \marrow \mathcal{G} \mathbb{R}$, but we require $I^\mathbb{N} \marrow \mathcal{G} \mathbb{R}$\footnote{actually, \lean{RV (⟪Ty.G Ty.real⟫)}, so we had omitted the finite
footprint property for understanding}. Note that in this case $\Delta = [] = $, an empty list so $\sem{\Delta} = \text{Unit}$ which is what \lean{List.TProd (⟪·⟫) []} is. We will have an empty context $\Delta$ for closed programs, which is
generally what we will see in a top level \texttt{wp} specification.
We compose \texttt{unif1.den} with the map \texttt{nil} to resolve the discrepancy.
But why do we bother with maps from $I^\mathbb{N}$ in the logic in the first place, when this is completely absent in the semantics of APPL? We refer the reader back to Sec \ref{sec:variables}, where this insightful question is answered.

\[
\{ \top \}\ \textsc{Unif1}\ \{ X.\ X \sim \mathrm{Unif}[0,1] \}
\]

The proof in Lilac using the Iris proof mode is as follows taken from \repo{Examples}:

\begin{leancode}

theorem unif1_spec :
    ⊢ wp (unif1.den ∘ᵣ nil) (λ X : RV ⟪Ty.real⟫ ↦ (X ∼ unif01_sem)) := by
  iapply wp_bind
  iapply wp_unif
  iintro %X HX
  iapply wp_ret
  rw [hrw]
  iexact HX

\end{leancode}

We present the proof state after each tactic\footnote{The interested reader is encouraged to clone the repository and view the file \repo{Examples} directly, in a Lean LSP enabled IDE such as vs-code.
Interactive proofs more informative when interacted with.}. Unlike the failed attempt in Sec \ref{sec:parametrised} where the proof state was artificially simplified for readability, here we faithfully reproduce the proof state
displayed by Lean's info-view, as seen by a user.

\texttt{iapply wp\_bind}
\begin{leancode}
⊢ 
⊢ wp (unif01.den ∘ᵣ nil) fun X ↦
  wp (#0.ret.den ∘ᵣ X ;; nil) fun X ↦ (X ∼ unif01_sem)
\end{leancode}

Application of the \texttt{wp\_bind} rule deconstructs the syntactic \texttt{bind} in the program \texttt{unif1} into nested \texttt{wp}
statements on the constituent programs \texttt{unif01} and \texttt{Term.ret \#0}. As expected the program of the inner \texttt{wp}
is not a closed-program and has the newly bound random variable $X$ in its context, at the head of the context. We can see already that
\texttt{\#0} is referring to $X$, but \texttt{\#0} is wrapped by \texttt{Term.ret}, so there is no reduction that can be applied yet.

\texttt{iapply wp\_unif}
\begin{leancode}
⊢ 
⊢ ∀ X, X ∼ unif01_sem -∗ 
  wp (#0.ret.den ∘ᵣ X ;; nil) fun X ↦ (X ∼ unif01_sem)
\end{leancode}

\texttt{iintro \%X HX}
\begin{leancode}
X : RV ⟪Ty.real⟫.carrier
⊢ 
∗HX : X ∼ unif01_sem
⊢ wp ((var head).ret.den ∘ᵣ X ;; nil) fun X ↦ (X ∼ unif01_sem)
\end{leancode}

Crucially here we are able to we have no \texttt{substLProp} unlike in Sec \ref{sec:parametrised}, unblocking
further application of rules. This is as far as we had previously gotten. 

\texttt{iapply wp\_ret}
\begin{leancode}
X : RV ⟪Ty.real⟫.carrier
⊢ 
∗HX : X ∼ unif01_sem
⊢ (#0.den ∘ᵣ X ;; nil) ∼ unif01_sem
\end{leancode}

We have removed the \texttt{Term.ret} allowing us to simplify the composition of program denotation and context.

\texttt{rw [hrw]}
\begin{leancode}
X : RV ⟪Ty.real⟫.carrier
⊢ 
∗HX : X ∼ unif01_sem
⊢ X ∼ unif01_sem
\end{leancode}

\texttt{hrw} is defined as
\begin{leancode}
lemma hrw {A : Ty} {X : RV ⟪A⟫} : ((#0).den ∘ᵣ (X ;; nil)) = X := by rfl
\end{leancode}

as seen the proof is \texttt{rfl} (reflexivity) which only holds for two terms which are \emph{definitionally equal},
meaning they reduce to the same normal form\footnote{there's quite a bit of nuance to what is considered a normal form
and which reductions need to be considered, for the purposes of computational efficiency of the \texttt{defEq} algorithm,
that we don't go into}. We can see a well defined procedure for constructing such rewrite lemmas, so this step can be automated
using a custom simplification tactic for our logic. However we didn't pursue this, as we would like to solve the deeper
problem of de Brujin indices being unintuitive, and custom tactics to automate rewrites would likely become redundant then.

\texttt{iexact HX}
\begin{leancode}
No goals
\end{leancode}

This concludes the first ever mechanised Lilac proof.

\section{Specifications with Expectations}
\label{sec:expectation}

Here we explore a more nuanced example where we first encountered difficulty with a notational side-stepping
of an issue in the Lilac paper~\citep{lilac}. Next, we this opportunity to also comment on a pervasive
simplification in proofs in our mechanisation and informal proofs provided in the paper. Here is a Hoare triple
style proof annotation of the program \texttt{Halve} which we have will next mechanise.

\[
\begin{array}{l}
\textcolor{blue}{\{\,\mathrm{own}\,X\,\}} \\[4pt]
Y \leftarrow \mathrm{unif}\ [0,1];\\
\textcolor{blue}{\{\,\mathrm{own}\,X * Y \sim \mathrm{Unif}[0,1]\,\}}\\
\mathrm{ret}\ (XY)\\
\textcolor{blue}{\{\,Z.\ \mathrm{own}\,X * Y \sim \mathrm{Unif}[0,1] * Z \circeq XY\,\}}\\
\textcolor{blue}{\{\,Z.\ \mathbb{E}[Y] = 1/2 \wedge \mathbb{E}[XY] = \mathbb{E}[X]\mathbb{E}[Y]\, * Z \circeq XY\,\}}\\
\textcolor{blue}{\{\,Z.\ \mathbb{E}[Y] = 1/2 \wedge \mathbb{E}[Z] = \mathbb{E}[X]\mathbb{E}[Y]\,\}}\\
\textcolor{blue}{\{\,Z.\ \mathbb{E}[Z] = \mathbb{E}[X]/2\,\}}
\end{array}
\]

We choose the \texttt{Halve} program, as not only does it exercise one of the salient features of the logic,
reasoning about expectations, it would also test whether the mechanisation can support this rather unconventional
goal of verifying a specification on an open variable. We are very much interested in whether our approach to variables
is robust to these kinds of examples, as it was a major difficulty that was overcome in this project.

The immediate difficulty to notice is that the expectation is a logical connective, $\mathbb{E}[X]=e$, includes $=$
in the syntax of the connective which can be misleading. It expects arguments typed as $E : \RV \mathbb{R}$ and $e : \mathbb{R}$.
So in particular, none of the expectations in the above proof, are using valid syntax. For example the post
condition must by expanded to:


$\textcolor{blue}{\{\,Z.\ \exists e,\ \mathbb{E}[Z] = e/2 \land \mathbb{E}[X] = e\,\}}$

% The verbosity of these shorthands unfolded quickly grows, with $\mathbbb{E}[XY] = \mathbb{E}[X] \mathbb{E}[Y]$ becoming

% $\exists e_X



% $\{\,Z.\ \exists e_Y e_x,\ \mathbb{E}[Y] = 1/2 \wedge \mathbb{E}[Z] = \mathbb{E}[X]\mathbb{E}[Y]\,\}$


Is this merely an inconvenience or a more fundamental problem, which might make this theorem unprovable in the logic.
The latter was in-fact a concern since it is not actually possible to instantiate the above existential! For clarity,
we emphasize that $\mathbb{E}[X]$ is not well-defined in the logic. But $\mathbb{E}[X]$ value is well-defined in the meta-logic (Lean), it is defined as an integral
of the random variable X. We want to instantiate $e$ with this semantic $\mathbb{E}[X]$. Let us first present the partial proof until we hit this roadblock.

\begin{leancode}
lemma hrw (X Y : RV ⟪Ty.real⟫) :
    (Term.den ((#1 * #0)) ∘ᵣ Y ;; X ;; nil) = X * Y := by rfl

abbrev post_half (X : RV ⟪Ty.real⟫) : LProp := wp (half.den ∘ᵣ (X ;; nil)) (λ Y : RV ℝ ↦
    iprop(∃ e, 𝔼[Y]=e/2 ∧ 𝔼[X]=e))

theorem half_spec (X : RV ⟪Ty.real⟫) : own X ⊢ post_half X := by
  iintro hX
  rw [post_half]
  iapply wp_bind
  iapply wp_unif
  iintro %Y hY
  iapply wp_ret
  rw [hrw]
\end{leancode}


but that is not possible as we demonstrate, with in this partia
i.e. an expression in Lean. This is well-defined. But we still aren't able to construc
We would of course want to instantiate it with the semantic (Lean )

\subsection{}

\section{Removing the almost surely equals connective ($\circeq$)}
As we have seen in the previous section we have completely removed the need for usage of the almost surely equals connective and analagous situations appear all throughout the paper,
both in examples using Core Lilac (Lilac without conditioning modality) and full Lilac. We give one more example from the paper to justify this claim, presenting an example using the
conditioning modality. Simultaneously, we are demonstrating that the set of definitions for Core Lilac will translate naturally into full Lilac.
We explain just what is relevant to showing that almost surely equals is not necessary, omitting the conditioning-specific proof rules. This is the first
instance in this report where we present something which does not have corresponding end-to-end proofs in Lean. Below is a condensed version of \citet[Appendix C]{lilac}:

\[
\begin{array}{l}
\textcolor{blue}{\{\, \top \,\}} \\[4pt]
Z \leftarrow \operatorname{flip}(1/2); \quad
X \leftarrow \operatorname{flip}(1/2); \quad
Y \leftarrow \operatorname{flip}(1/2); \\
\textcolor{blue}{\{\, Z \sim \operatorname{Ber}(1/2) \ast X \sim \operatorname{Ber}(1/2) \ast Y \sim \operatorname{Ber}(1/2) \,\}} \\[4pt]
A \leftarrow \mathrm{ret}\, (X \lor Z); \quad
B \leftarrow \mathrm{ret}\, (Y \lor Z); \\
\textcolor{blue}{\{\, Z \sim \operatorname{Ber}(1/2) \ast X \sim \operatorname{Ber}(1/2) \ast Y \sim \operatorname{Ber}(1/2) \ast A = (X \lor Z) \ast B = (Y \lor Z) \,\}} \\[4pt]
\operatorname{ret}\, (Z,X,Y,A,B) \\
\textcolor{blue}{\{\, \underset{z \leftarrow Z}{\mathbf{C}}(\operatorname{own}(A) \ast \operatorname{own}(B)) \,\}}
\end{array}
\]

We translate the Hoare style specific to a \texttt{wp} style specification, and show the stages of reduction as we apply the \texttt{wp} proof rules:

\[
\begin{array}{l}
\textcolor{blue}{\vdash \mathrm{wp}} \\
\left(
\begin{array}{l}
Z \leftarrow \operatorname{flip}(1/2); \quad
X \leftarrow \operatorname{flip}(1/2); \quad
Y \leftarrow \operatorname{flip}(1/2); \\
A \leftarrow \mathrm{ret}\,(X \lor Z); \quad
B \leftarrow \mathrm{ret}\,(Y \lor Z); \\
\operatorname{ret}\, (Z,X,Y,A,B)
\end{array}
\right) \\[4pt]
\textcolor{blue}{\{\, (Z,X,Y,A,B).\ \underset{z \leftarrow Z}{\mathbf{C}}(\operatorname{own}(A) \ast \operatorname{own}(B)) \,\}}
\end{array}
\]

\[
\begin{array}{l}
\textcolor{blue}{X\, Y\, Z : \RV \mathbb{R}} \\
\textcolor{blue}{X \sim Ber (1/2)} \\
\textcolor{blue}{Y \sim Ber (1/2) }\\ 
\textcolor{blue}{Z \sim Ber (1/2)} \\
\textcolor{blue}{\vdash \mathrm{wp}} \\
\left(
\begin{array}{l}
A \leftarrow \mathrm{ret}\,(X \lor Z); \quad
B \leftarrow \mathrm{ret}\,(Y \lor Z); \\
\operatorname{ret}\, (Z,X,Y,A,B)
\end{array}
\right) \\[4pt]
\textcolor{blue}{\{\, (Z,X,Y,A,B).\ \underset{z \leftarrow Z}{\mathbf{C}}(\operatorname{own}(A) \ast \operatorname{own}(B)) \,\}}
\end{array}
\]

\[
\begin{array}{l}
\textcolor{blue}{X\, Y\, Z : \RV \mathbb{R}} \\
\textcolor{blue}{X \sim Ber (1/2)} \\
\textcolor{blue}{Y \sim Ber (1/2) }\\ 
\textcolor{blue}{Z \sim Ber (1/2)} \\
\textcolor{blue}{\vdash \mathrm{wp}}
\ \operatorname{ret}\, (Z,X,Y,(X \lor Z),(Y \lor Z)) \\
\textcolor{blue}{\{\, (Z,X,Y,A,B).\ \underset{z \leftarrow Z}{\mathbf{C}}(\operatorname{own}(A) \ast \operatorname{own}(B)) \,\}}
\end{array}
\]

\[
\begin{array}{l}
\textcolor{blue}{X\, Y\, Z : \RV \mathbb{R}} \\
\textcolor{blue}{X \sim Ber (1/2)} \\
\textcolor{blue}{Y \sim Ber (1/2)} \\
\textcolor{blue}{Z \sim Ber (1/2)} \\
\textcolor{blue}{\vdash \{\, \underset{z \leftarrow Z}{\mathbf{C}}(\operatorname{own}(X \lor Z) \ast \operatorname{own}(Y \lor Z)) \,\}}
\end{array}
\]

\Todo{Condense into a diagram with arrows stating the proof rules applying to achieve the transformations}

We have made the program context (of the form \lean{∘ᵣ (... ;; nil)}) implicit, and given explicit
names rather than de Brujin indices to variables, as a consequence of this omission, for readability. We color code
the program in black and the LProp logical assertions in \textcolor{blue}{blue} to show the entire reduction of
of the program using weakest precondition reasoning and our proof rules, never needs to use $\circeq$ the almost
surely equals connective. The proof is not completed, and we have stopped at the point that we would have needed proof
rules about the conditioning modality, the mechanisation of whose soundness proofs, we leave as future work\footnote{This
is expected to be similar in difficulty to the formalisation of \texttt{wp\_unif} possibly requiring the mechanisation
of relevant measure theoretic lemmas first before we can proceed with the soundness proof itself.
}

\subsection{Comparison with Bluebell}
\Todo{}
This section is incomplete. What I aim to talk about is that the definition of almost surely equals in Lilac,
and as a consequence the proof its properties such as being an equivalence relation, got somewhat convoluted,
and in Bluebell a fundamentally alternative approach was pursued in part to offer a cleaner solution to this problem.

In particular Bluebell includes variables in its resource algebra. @Shing Hin can you confirm if these claims are correct?

Because of this observation, I just want to state that it is strange I ended up not needed almost surely equals at all.
And I would want to try out more example and implement full Lilac (with conditioning) before being confirming if
almost surely equals is indeed not necessary....